\section{Technology stack and justification}
Table~\ref{tab:stack} lists every major tool, why it fits this use case and what we compared it with. We
chose each tool for a need of the fleet problem, not because it is popular. All versions are the current
stable releases in September 2026 \cite{spark,kafka,airflow}.

\begin{table}[!htb]
\centering\small
\caption{Technology stack.}
\label{tab:stack}
\begin{tabularx}{\textwidth}{@{}P{1.9cm} P{2.5cm} L P{2.8cm}@{}}
\toprule
\textbf{Layer} & \textbf{Choice} & \textbf{Why it fits our use case} & \textbf{Compared with} \\
\midrule
Ingestion & Apache Kafka 4.3 (KRaft, native image) &
A durable, ordered log. Keys keep each vehicle's events in order, and consumers can re-read by offset after
a failure, so either side can restart without losing data. KRaft needs no ZooKeeper. &
RabbitMQ (no replay by offset), Redpanda \\
Stream processing & Apache Spark 4.1 Structured Streaming &
Event-time windows and watermarks for late GPS data, de-duplication, stream-static joins and checkpoints
for exactly-once processing. The same API runs the batch layer. &
Apache Storm (no batch API, harder state handling) \\
Batch processing & Apache Spark 4.1 (batch) &
Recomputes a day from Parquet with window functions (time in each status, idle streaks) and joins the cost
file, with the module shared with the speed layer. &
Plain SQL in PostgreSQL (raw events would fill the serving store) \\
Orchestration & Apache Airflow 3.3 &
Sensors wait for the file, tasks retry, failures raise an alert, and any day can be rerun with one
parameter. The web UI shows every run. &
cron (no dependencies, retries or run history) \\
Master dataset & Parquet files on a Docker volume &
Columnar, compressed and partitioned by day, so it is cheap to keep forever. In production the same layout
would sit on S3 or HDFS. &
Keeping all history in Kafka \\
Serving store & PostgreSQL 17 &
Results are small (hundreds of rows a day). We need joins for the merge views, upserts for streaming updates
and transactions to publish a whole day at once. &
Cassandra (no joins or multi-row transactions) \\
Serving & FastAPI, Grafana 12, HTML report &
A documented REST API with \code{/metrics}, live dashboards on PostgreSQL and Prometheus, and a report that
is easy to email to managers. &
Streamlit \\
Observability & Prometheus 3, Alertmanager, StatsD exporter, JSON logs &
Metrics from every stage and alert rules. Alertmanager groups and routes alerts, and StatsD brings in the
Airflow metrics. &
Loki and Jaeger (too much memory for our laptop) \\
Deployment & Docker Compose &
One command starts 13 services (11 long-running and 2 set-up jobs) with health checks and memory limits. &
Kubernetes (too heavy for one laptop) \\
\bottomrule
\end{tabularx}
\end{table}

\textbf{Designing for an 8 GB laptop.} Our computer has 8\,GB of memory and Docker gets only about
3.8\,GB. This constraint shaped some choices. We run the three Airflow components in one container with the
LocalExecutor instead of the official set-up with Celery and Redis. Kafka uses the official GraalVM
\emph{native} image, which needed less than 150\,MB in our run instead of about 350\,MB for the Java image. Both Spark jobs
run in local mode with a 512\,MB heap and use only the fast C1 compiler of the Java virtual machine. Without
these changes, the system swapped memory and one batch run took eight minutes.
